\chapter[Conclusion]{Conclusion and Open Questions}
\label{chapter:conclusion}

Algorithmic information theory continues to offer insights into geometric measure theory. In this dissertation, we have reviewed a subset of the existing results, methods, and notions relevant for witnessing this connection. We have proved some extra robustness and relations in the effective framework, and have extended AIT to a broader class of metric spaces. Namely, net spaces are a productive setting over which to work with algorithmic information theory. And effective notions most closely simulate their classical counterparts over net spaces which are sufficiently rich with net measures. This is exemplified by the various point-to-set principles. 

\section{Dependence on the Net or Dense Subset}

Certain notions in algorithmic information theory explicitly depend on the choice of a computable dense subset (or net) with respect to which approximations are taken. For instance, as introduced in Section~\ref{section:Lifting_Prefix_Complexity_to_Reals}, the standard lift of prefix complexity to Euclidean space bases an arbitrary subset's complexity solely on the rational points it contains. Similarly, N.~Lutz's locally-optimal outer measure $\boldsymbol{\kappa}(X) := 2^{-K(X)}$ satisfies $\boldsymbol{\kappa}(X) = \boldsymbol{\kappa}(X \cap \mathbb{Q}^m)$, as well as some finiteness properties specific to $\mathbb{Q}^m$. How would another choice of computable dense subset affect these notions of complexity and measure? Lemma~4.8 in \cite{LL20} makes clear that multiplicative domination (in the sense of Definition~\ref{definition:domination_over_LDS}) on dyadic cubes is equivalent to multiplicative domination on the collection of open balls of dyadic rational radius; and their Theorem~4.5 implies that their choice of $\boldsymbol{\kappa}$ based on $\mathbb{Q}^m$ dominates on both these nets. But for other choices of computable dense subset, it is not clear that the corresponding lift of prefix complexity would assign comparable complexities to each subset of $\mathbb{R}^m$, nor comparable measures to each subset according to the corresponding outer measure $\boldsymbol{\kappa}$. Thus, the local dimension notion associated to $\boldsymbol{\kappa}$ would not necessarily match with other effective dimension notions. Analogous questions may be asked over generic net spaces.  

\section{Conditional, Mutual Dimension}

In Section~\ref{section:Conditional_and_Mutual_Dimensions}, we reviewed two variants of effective dimension: \textit{mutual} and \textit{conditional} dimensions. Notice that for finitary objects, it was indeed possible to define a \textit{conditional, mutual information} of the form $I(a : b \mid c)$: this intuitively measures the information content of the pair of objects $a$ and $b$ when given the data $c$. But mutual dimension is only defined in the unconditional sense. Can a useful notion of conditional, mutual dimension be defined? If $m, n, \ell \in \omega$, $x \in \mathbb{R}^m$, $y \in \mathbb{R}^n$, $z \in \mathbb{R}^\ell$, and $r, s, t \in \omega$, then one reasonable candidate definition would be:
\begin{align*}
    \operatorname{mdim}(x : y \mid z) &:= \liminf_{r \to \infty} \frac{I_{r : r \mid r}(x : y \mid z)}{r}, \quad \quad I_{r : s \mid t}(x : y | z) := \min_u \seq{\min_{p, q} \seq{I(p : q \mid u)}},
\end{align*}
where the choices from $p$, $q$, and $u$ range in: $p \in B_{2^{-r}}(x) \cap \mathbb{Q}^m$, $q \in B_{2^{-s}}(y) \cap \mathbb{Q}^n$, and $u \in B_{2^{-t}}(z) \cap \mathbb{Q}^\ell$. The definition for $I_{r : s \mid t}$ is analogous to that of the conditional prefix complexity $K_{r \mid s}$. 

Recall that conditional complexity admits an approximation by $K$-minimizers. That is, $K_{r \mid s}(x \mid y) \approx K(p^* \mid q^*)$, where $p^*$ and $q^*$ are $K$-minimizers of the balls $B_{2^{-r}}(x)$ and $B_{2^{-s}}(y)$, respectively (see Proposition~\ref{lemma:conditional_K_minimizers_suffice}). But in the case of conditional, mutual information, there is not as clear of a path towards showing the analogous statement: $I_{r : s \mid t}(x : y \mid z) \stackrel{?}{\approx} I(p^* : q^* \mid w^*)$, where each $p^*, q^*$, and $w^*$ are $K$-minimizers of the appropriate balls about $x$, $y$, and $z$, respectively. To be specific, while a quick calculation shows the approximate inequality $I_{r : s \mid t}(x : y \mid z) \geq I(p^* : q^* \mid w^*)$ in one direction, we are missing a proof of the reverse direction. The main issue is the lack of known results analogous to the linear sensitivities of conditional complexity: Lemmas~\ref{minmax_linear_sensitivity_in_s} and~\ref{minmax_linear_sensitivity_in_r}.

If the approximation to $I_{r : s \mid t}$ by $K$-minimizers does hold up to logarithmic terms, then an analogous result to Theorem~\ref{theorem:SOI_conditional_complexity_r_s_t} should follow for mutual information, and one could go on to define a \textit{conditional, mutual dimension} and compare it to the other variants to effective dimension.

\section{Applications of CALFs}

Recall from Chapter~\ref{chapter:CALFs} that a continuous absolutely Lipschitz family (CALF) generalizes the uniform continuity properties satisfied by the family of non-vertical planar lines expressed in slope-intercept form. Generalizing the results of N. Lutz and D. Stull for planar lines, we have shown that every CALF admits both a finitary theorem (\ref{theorem:finitary}) on the Kolmogorov complexity of points along its graph, as well as an infinitary theorem (\ref{theorem:infinitary}) on the effective dimension of points along its graph. 

In \cite{LS16}, N.~Lutz and D.~Stull used their infinitary theorem (Theorem~1.1) and the Point-to-Set Principle~\ref{theorem:PTS_for_H_dimension} to demonstrate a new lower-bound for the Hausdorff dimension of generalized Furstenberg sets. Given $d, \delta \in [0, 1]$, the class of generalized Furstenberg sets for $d$ and $\delta$ is denoted $F_{d, \delta}$, and a subset $X \subseteq \mathbb{R}^2$ qualifies for $X \in F_{d, \delta}$ if there is some $J \subseteq \mathbb{S}^1$ such that $\dim_{\operatorname{H}} J \geq d$ and such that for any $e \in J$, there exists a line $l_e$ in the direction $e$ for which $\dim_{\operatorname{H}}(X \cap l_e) \geq \delta$. Then given $d, \delta > 0$ and $X \in F_{d, \delta}$, Theorem~4.3 in \cite{LS16} states $\dim_{\operatorname{H}} X \geq \delta + \min \seq{d, \delta}$.

For a more general family of maps such as a CALF $\Phi$ on a domain $\Omega \times \Xi \subseteq \mathbb{R}^m \times \mathbb{R}^\ell$, one could similarly define a collection of subsets $F^\Phi_{d, \delta}$ for any $0 \leq d \leq m$ and $0 \leq \delta \leq \ell$ as follows: $X \subseteq \mathbb{R}^n$ satisfies $X \in F^\Phi_{d, \delta}$ whenever there exists $J \subseteq \Omega$ with $\dim_{\operatorname{H}}J \geq d$ and for all $\boldsymbol{\alpha} \in J$, there exists another $\boldsymbol{\alpha}' \in \Omega$ with $\dim(\boldsymbol{\alpha}') = \dim(\boldsymbol{\alpha})$ such that $\dim_{\operatorname{H}}(X \cap \Phi^{\boldsymbol{\alpha}'}) \geq \delta$. It is not clear for a given CALF $\Phi$ whether $F^\Phi_{d, \delta}$ will be non-empty, or whether $F^\Phi_{d, \delta}$ will have a useful interpretation in the context of geometric measure theory. But the proof method of Theorem~4.3 in \cite{LS16} likely extends for such sets for non-empty $F^\Phi_{d, \delta}$ when $d, \delta > 0$, and should make use of Theorem~\ref{theorem:infinitary}.

Each new example of a CALF provides a new geometry over which to consider possible implications for classical geometric measure theory. Certain families of planar isometries are natural candidates for CALFs. While the scaling Lipschitz continuity property is more or less clear for families of isometries, we run into difficulties demonstrating the other conditions. For instance, proving the scaling co-Lipschitz continuous differences property for the family of planar rotations is made difficult by the situation that the difference between two angle parameters can be large while the angular difference between them is small. Additionally, it is not clear that one may parameterize the family of planar reflections in a $1$-dimensional subspace on an open domain fit for describing them by a CALF. 

It is possible that certain results from classical geometric measure theory admit corresponding finitary results. That is, we imagine any statement about the Hausdorff dimension of sets in Euclidean space having a corresponding statement about the Kolmogorov complexity of finitary inputs and outputs as in Theorem~\ref{theorem:finitary}. And that finitary result should imply a similar statement but for effective dimension as in Theorem~\ref{theorem:infinitary}. And by applying the Point-to-Set Principle, one could recover the classical result for Hausdorff dimension. We hope to see more examples of this finitary-infinitary structure for more results in GMT. 

\section{Effectivity for the Subsets of Exact Measure}

As discussed in Section~\ref{section:subsets_of_exact_measure}, we have extened Besicovitch's result about finding compact subsets of finite, non-zero $s$-measure to sufficiently nice nets in our Theorem~\ref{theorem:existence_of_compact_subset_exact_for_premeasure}. The proof of this result relies on a density result: Proposition~\ref{proposition:full_density_of_U_and_S}. 

The \textit{$n$-thin} sequences defined in our proof play the crucial role of helping pass from $\mathcal{U}^c_n$ to $\mathcal{U}^c_{n _ 1}$, as shown in Lemma~\ref{lemma:use_n_thin_to_pass_n_to_n+1}. They have the special property that any length-$n$ string used in a cover of $E_Z$ may be replaced by proper extensions to make another cover of $E_Z$ not significantly greater in its direct $\rho_h$-weight. The $n$-thin sequences were shown to be dense in any non-empty $\mathcal{S}^d_{n_0}$ whenever $n \geq n_0$ in Lemma~\ref{lemma:n_thins_dense_in_complement}. We suspect that the $n$-thin sequences should be useful in other contexts involving making refinements to optimal covers. 

One could also imagine effectivizing the proof of Theorem~\ref{theorem:existence_of_compact_subset_exact_for_premeasure} as follows: finding an oracle $B \in 2^\omega$ such that, if $F \subseteq 2^\omega$ were in fact a $\Pi^0_1$-class, $B$ is strong enough to compute a code $Z \in \mathcal{P}_F$ for which $E_Z$ satisfies the desired properties, putting $E_Z \in \Pi^0_1(B)$. One might even modify $\phi_{n, m}(Z)$ to instead optimize over the \textit{a priori} discrepancy $\operatorname{\Delta M}_s(\sigma)$ over the strings for which $Z(\sigma) = 1$. Then, an appropriate lightface version (i.e., correspondence principle) of Lutz and Miller's point-to-set principle~\ref{theorem:point_to_set_s_measure} would relate this function back to the Hausdorff $s$-measure of the compact subset $E_Z$ coded by $Z$. In particular, we hope for an effective proof of Besicovitch's result over Cantor space via lower-semicomputable continuous semimeasures. 